% Generated by scripts/assemble.py; corrections belong in source/corrections.json.
\examhead{5}{北京交通大学 2024 年数学分析试卷(回忆版)}{12}
% p013-b05; source PDF page 13
\par\Needspace{5\baselineskip}
\blockanchor{p013-b05}
\majorquestion{一}数列 $\{a_{n}\}$ 没有最大值, 没有最小值, 证明: $\{a_{n}\}$ 发散.\par
% p013-b06; source PDF page 13
\par\Needspace{5\baselineskip}
\blockanchor{p013-b06}
\majorquestion{二}$f(x)$ 在 $[0,2]$ 上连续，且 $f(0) + f(1) = 2, f(2) = 1$. 证明：存在 $\xi \in (0,2)$ 使得 $f'(\xi) = 0$.
\majorquestion{三}计算 $\int \frac{x\mathrm{e}^x}{(x + 1)^2}\mathrm{d}x,\int \frac{1}{(x - 1)^2(x - 2)}\mathrm{d}x.$\par
% p013-b07; source PDF page 13
\par\Needspace{5\baselineskip}
\par\Needspace{6\baselineskip}
\blockanchor{p013-b07}
\majorquestion{四}设 $f(x)$ 在 $[a, b]$ 上三阶可导, 证明存在 $\xi \in (a, b)$ 使得\par
\nopagebreak[4]
% p013-b08; source PDF page 13
\blockanchor{p013-b08}
\displayformula{f (b) = f (a) + \frac {1}{2} (b - a) \left[ f ^ {\prime} (a) + f ^ {\prime} (b) \right] - \frac {1}{12} (b - a) ^ {3} f ^ {\prime \prime \prime} (\xi).}
% p013-b09; source PDF page 13
\par\Needspace{5\baselineskip}
\blockanchor{p013-b09}
\majorquestion{五}证明:\par
% p013-b10; source PDF page 13
\blockanchor{p013-b10}
\question{1}$f(x) = \frac{|\sin x|}{x}$ 在 $(-1,0)$ 和 $(0,1)$ 上一致连续.\par
% p013-b11; source PDF page 13
\blockanchor{p013-b11}
\question{2}$f(x) = \frac{|\sin x|}{x}$ 在 $(-1,0)\cup (0,1)$ 上不一致连续.\par
% p013-b12; source PDF page 13
\par\Needspace{5\baselineskip}
\par\Needspace{6\baselineskip}
\blockanchor{p013-b12}
\majorquestion{六}设\par
\nopagebreak[4]
% p013-b13; source PDF page 13
\blockanchor{p013-b13}
\displayformula{f (x) = \left\{ \begin{array}{l l} x y \sin \frac {1}{x ^ {2} + y ^ {2}}, & x ^ {2} + y ^ {2} \neq 0; \\ 0, & x ^ {2} + y ^ {2} = 0. \end{array} \right.}
% p013-b14; source PDF page 13
\blockanchor{p013-b14}
\question{1}求  $\frac{\partial f}{\partial x}, \frac{\partial f}{\partial y}.$\par
% p013-b15; source PDF page 13
\blockanchor{p013-b15}
\question{2}$\frac{\partial f}{\partial x}, \frac{\partial f}{\partial y}$ 在 $(0,0)$ 是否连续, $f(x,y)$ 在 $(0,0)$ 是否可微?\par
% p013-b16; source PDF page 13
\par\Needspace{5\baselineskip}
\blockanchor{p013-b16}
\majorquestion{七}证明 $\left\{ \begin{array}{l} x + y + u + v = 0; \\ x^2 + y^2 + u^4 + v^4 = 1. \end{array} \right.$ 存在 $\left( \begin{array}{l} u \\ v \end{array} \right) = \left( \begin{array}{l} u(x, y) \\ v(x, y) \end{array} \right)$, 并计算 $\frac{\partial u}{\partial x}, \frac{\partial^2 u}{\partial x^2}$.\par
% p013-b17; source PDF page 13
\par\Needspace{5\baselineskip}
\blockanchor{p013-b17}
\majorquestion{八}求三重积分 $\iiint_{\Omega} (x + y + z) \, \mathrm{d}x \, \mathrm{d}y \, \mathrm{d}z,$ 其中 $\Omega$ 为(具体数据忘了).\par
% p013-b18; source PDF page 13
\par\Needspace{5\baselineskip}
\blockanchor{p013-b18}
\majorquestion{九}求曲线积分问题, 应该是考察斯托克斯公式的应用.\par
% p013-b19; source PDF page 13
\par\Needspace{5\baselineskip}
\blockanchor{p013-b19}
\majorquestion{十}设 $f_{n}(x) = \frac{nx}{1 + (1 + n^{2})x^{2}}$.\par
% p013-b20; source PDF page 13
\blockanchor{p013-b20}
\question{1}证明: $\{f_n(x)\}$ 在 $(0,1]$ 上不一致收敛.\par
% p013-b21; source PDF page 13
\blockanchor{p013-b21}
\question{2}证明: $\{f_n(x)\}$ 在 $[\alpha, 1]$ 上一致收敛, 其中 $0 < \alpha < 1$.\par
% p013-b22; source PDF page 13
\par\Needspace{5\baselineskip}
\blockanchor{p013-b22}
\majorquestion{十一}计算 $\sum_{n=1}^{\infty} \frac{1}{(n+1)2^n}$.\par
